\section{Related Work}
\label{sec:rw}

This section surveys the literature across five themes: foundational attention mechanisms and transformers (Section~\ref{subsec:attention}), time-series forecasting with neural networks (Section~\ref{subsec:timeseries}), cloud and cluster monitoring (Section~\ref{subsec:monitoring}), SLA violation prediction (Section~\ref{subsec:sla}), and proactive resource management (Section~\ref{subsec:proactive}). The baseline paper for this project (Section~\ref{subsec:baseline}) is then positioned within this landscape, and the specific gap this work addresses is articulated (Section~\ref{subsec:gap}).

\subsection{Attention Mechanisms and Transformers}
\label{subsec:attention}

The transformer architecture, introduced by Vaswani et al.~\cite{Vaswani17} in ``Attention Is All You Need,'' replaced recurrent and convolutional layers with multi-head self-attention mechanisms that compute pairwise relationships between all tokens in a sequence. This allows the model to capture long-range dependencies without the vanishing gradient problems that plague recurrent neural networks. The multi-head design partitions the model's representational capacity into multiple parallel attention ``heads,'' each of which can specialize on different aspects of the input. This architecture has become the foundation for modern natural language processing (e.g., BERT, GPT) and has increasingly been applied beyond text to domains such as time-series forecasting, computer vision, and system monitoring.

Multi-head attention has proven particularly effective for time-series forecasting. MSGNet~\cite{MSGNet24} applies multi-head self-attention to learn multi-scale inter-series correlations in multivariate time series, using Fast Fourier Transformation (FFT) to project data into scale-specific spaces and a learnable adjacency matrix per time scale. Wu et al.~\cite{Wu25} revisit the conventional attention mechanism for multivariate time-series forecasting (MTSF), proposing Frequency Spectrum attention (FSatten) and Scaled Orthogonal attention (SOatten), both of which outperform conventional attention without requiring architectural changes. MixNet~\cite{MixNet26} introduces a scale-adaptive multi-head attention mechanism based on a mixture-of-experts framework, dynamically assessing the significance of each attention module for flexible feature extraction across diverse time-series data. Liu et al.~\cite{BilinearAttention22} propose a multi-head temporal attention-augmented bilinear network for financial time-series prediction, demonstrating that simultaneous focus on multiple temporal instances enhances prediction performance.

Despite these successes, a recurring limitation in many multi-head architectures is that heads operate independently: once the input is projected into separate head subspaces, those heads do not exchange information with one another. This project directly addresses that limitation.

\subsection{Time-Series Forecasting with Neural Networks}
\label{subsec:timeseries}

Time-series prediction has long been central to monitoring and forecasting tasks in distributed systems. Early neural approaches include Elman's simple recurrent networks~\cite{Elman90}, which introduced the idea of using recurrent connections to maintain a dynamic memory of prior hidden states, enabling the network to represent temporal structure implicitly. However, these vanilla recurrent networks suffer from the vanishing gradient problem when learning long-term dependencies.

Hochreiter and Schmidhuber~\cite{Hochreiter97} addressed this with Long Short-Term Memory (LSTM) networks, which use gating mechanisms (input, forget, and output gates) to regulate information flow and maintain constant error flow over long sequences. LSTMs have since become a workhorse for time-series forecasting in cloud systems. For example, the workload prediction model by Balraj and Chandran~\cite{WorkloadmLSTM24} uses multiplicative LSTM (mLSTM) to predict future workloads in cloud data centers, addressing long-term dependencies to reduce SLA violations. Similarly, Alaa'anzy and Othman~\cite{MultiModelDL19} evaluate GRU, LSTM, and Bi-directional LSTM for cloud resource usage prediction to support proactive workflow adaptation, selecting the most accurate model per task based on runtime performance.

While recurrent architectures like LSTM excel at capturing temporal dependencies, they process sequences sequentially, which limits parallelization and can be slow for long sequences. Transformers, by contrast, process all timesteps in parallel via self-attention, offering both computational efficiency and the ability to capture long-range dependencies directly. This has driven recent interest in transformer-based models for time-series tasks, including the multi-head attention forecasting approaches discussed in Section~\ref{subsec:attention}.

\subsection{Cloud and Cluster Monitoring}
\label{subsec:monitoring}

Modern cloud clusters generate vast volumes of telemetry data (CPU, memory, disk I/O, network throughput) from thousands of nodes, requiring scalable, automated monitoring systems. Traditional approaches rely on fixed thresholds applied to individual metrics, raising alerts only after a metric has breached its limit -- a reactive strategy that provides no lead time for intervention.

The Google Borg cluster traces~\cite{Tirmazi20,Liu12} have become a foundational benchmark for workload characterization and scheduler evaluation. The 2011 trace provided 29 days of job submissions, scheduling decisions, and resource usage from a single 12.5k-machine Borg cell; the 2019 trace extended this to eight cells (~96k machines) with richer metadata including allocation sets, job dependencies, and vertical scaling events. Tirmazi et al.~\cite{Tirmazi20} report that the 2019 workload shows a 3.7× higher job submission rate and 7× increase in tasks needing scheduling compared to 2011, despite comparable cell sizes, underscoring the growing complexity of modern cluster management.

Container orchestration platforms like Kubernetes have driven demand for observable, metrics-driven autoscaling. Prometheus, an open-source time-series database and monitoring system, has become the de facto standard for Kubernetes telemetry collection. Kotsur and Yushchenko~\cite{PrometheusKubernetes20} propose an automatic anomaly detection system for Kubernetes using Prometheus metrics, demonstrating that machine learning over collected metrics can detect anomalies earlier than human operators. Li et al.~\cite{KubernetesHPA20} evaluate Kubernetes Horizontal Pod Autoscaler (HPA) with Prometheus custom metrics, showing that flexible metric selection improves autoscaling responsiveness. Makarov et al.~\cite{ProactiveKubernetes24} introduce a proactive horizontal scaling method for Kubernetes that forecasts future CPU usage from Prometheus time series using Prophet and ARIMA, scaling resources before demand spikes occur. Huang and Pierre~\cite{KubernetesFederation24} address aggregate monitoring for geo-distributed Kubernetes cluster federations, proposing metrics deduplication and aggregation strategies to reduce monitoring overhead.

Anomaly detection in cloud telemetry has also advanced through spatiotemporal learning. Qiao et al.~\cite{ADDSTL23} propose AD-DSTL, which combines graph convolutional networks (GCNs) to model service dependencies with LSTMs to capture temporal dynamics, achieving high recall and precision (~0.9) on multiple cloud system datasets. Scheinert et al.~\cite{CloudAnomalyIBM24} present a large-scale anomaly detection dataset from IBM Cloud (39,365 five-minute intervals, 117,448 features), evaluating GRU-based autoencoders that detect anomalies up to 20 minutes ahead of traditional threshold-based monitoring.

\subsection{SLA Violation Prediction}
\label{subsec:sla}

Service Level Agreements (SLAs) define contractual performance guarantees (e.g., uptime, latency, resource availability) with penalty clauses for violations. Predicting SLA violations before they occur enables cloud providers to take corrective action, avoiding penalties and maintaining customer trust.

Early work applied traditional machine learning to SLA prediction. Ivan et al.~\cite{SLAViolationML19} compare Hidden Markov Models, Neural Networks, K-Nearest Neighbors, and Support Vector Machines for workload prediction to estimate SLA violations, demonstrating that multiple learning techniques can predict violations with reasonable confidence. Kumar and Selvakumar~\cite{SLAViolationSCG19} use a scaled conjugate gradient neural network with oversampling to handle biased datasets, achieving 96.76\% and 93.23\% classification accuracy on two real-world SLA violation datasets. Hani et al.~\cite{ManifoldSLA17} apply manifold learning to reduce high-dimensional QoS parameters to a 1-D violation level (low/medium/high), using Support Vector Regression for time-series prediction with 80\% accuracy on 14 days of QoS data.

More recent approaches leverage deep learning and graph-based models. Karras et al.~\cite{GNNSLA23} propose a Graph Neural Network (GNN)-based model that captures client associativity (context representation) alongside temporal QoS metrics, demonstrating that graph-based models significantly improve SLA violation prediction accuracy over traditional vector-based and sequential models. Balraj and Chandran~\cite{WorkloadmLSTM24} use multiplicative LSTM (mLSTM) to predict workload and reduce SLA violations, outperforming other LSTM variants on long-term dependency tasks.

\subsection{Proactive Resource Management}
\label{subsec:proactive}

Proactive resource management anticipates future demand and provisions resources ahead of time, contrasting with reactive approaches that respond only after thresholds are breached. This requires accurate workload forecasting and automated scaling mechanisms.

Navarro et al.~\cite{SVMProvisioning19} propose a predictive auto-scaling mechanism based on Support Vector Machine (SVM) regression combined with queue theory, demonstrating that SVM handles both linear and nonlinear workload patterns and avoids the local minima issues of neural networks. Sun et al.~\cite{ProactiveAllocation21} introduce a runs-test-based adaptive prediction algorithm (EEMD-ARIMA) that selects the most accurate forecasting method via an adaptive strategy, allocating resources proactively for sudden requests to guarantee timeliness. Naviner et al.~\cite{ProactiveMicroservices24} design a proactive management system for microservices that autonomously predicts load variations and resource scarcity, making resources available to ensure continuity without interruptions in AWS environments. Aldossary and Sotiriadis~\cite{AdaptiveAutoscaling23} propose an adaptive auto-scaling framework (ADA-RP) that clusters workloads into High/Medium/Low demand using K-means, trains a CNN model to predict future demand, and applies dynamic time warping to generate real-time adaptive scaling plans without introducing booting delays.

Alaa'anzy and Othman~\cite{MultiModelDL19} evaluate multiple deep learning models (GRU, LSTM, Bi-directional LSTM) for cloud resource usage prediction to support proactive workflow adaptation, selecting the model with the highest accuracy (lowest RMSE, MAE, MAPE) per task at runtime. These works collectively demonstrate that proactive, predictive approaches consistently outperform reactive threshold-based monitoring in terms of lead time, resource efficiency, and SLA compliance.

\subsection{Baseline: Multi-Head Attention for Telemetry-Based Violation Prediction}
\label{subsec:baseline}

The baseline family committed by formal CA2 is hybrid deep learning and traditional monitoring for cloud task/cluster failure prediction -- notably Aldomi et al.~\cite{Aldomi26} (SelectKBest feature selection + GRU temporal extraction + ML classifiers on Google Cluster Trace) and related RF/KNN/SVM/GRU monitors. Thapliyal~\cite{Thapliyal26} (per-metric multi-head SLA monitoring) is retained only as \textbf{related architecture} for the MHSA encoder; it is \textbf{not} the binding CA2 baseline.

This project's formal evaluation target remains GCT-backed Acc/Prec/Rec/F1/ROC-AUC/latency against that hybrid/traditional family. A 2011 subset is scored in Section~\ref{sec:Evaluation}; 2019 Borg cells used by Aldomi et al.\ are not in this repository.

The paper encodes SLA rules as structured JSON objects to generate training labels without manual annotation, trains a multi-head transformer model where each head specializes in one SLA rule, and deploys the trained model as a streaming inference service that emits structured prediction events. These events are transformed into three role-specific views: finance schemas exposing credit liability, operations schemas surfacing risk scores and recommended interventions, and compliance schemas bundling predictions with immutable telemetry signatures for audit.

Critically, the paper's own reported results expose a specific weakness of the strict one-head-per-metric design: the power-draw head shows ``a skewed, bimodal error profile ... with systematic underprediction during high-load transients, a direct consequence of instantaneous load volatility.'' The authors attribute this to the architecture's decoupled heads never exchanging information with each other, even though their own data shows power and temperature are correlated. This is the precise gap this project targets.

\subsection{The Gap This Project Targets}
\label{subsec:gap}

Because each head in a strict per-metric design can be trained and read out independently, a head may under-use another metric's precursor signal. Thapliyal~\cite{Thapliyal26} reports such underprediction on SLA telemetry; the artefact optionally adds a cross-head fusion layer as an architectural experiment on synthetic data. That experiment is orthogonal to -- and does not satisfy -- the formal CA2 requirement to score MHSA-TDL on Google Cluster Trace against Aldomi-style and classical monitors.

Traditional reactive threshold monitoring is retained as a non-learned monitor. Classical RF/KNN/SVM and an Aldomi SelectKBest+GRU+RF path are scored on the same GCT splits. The Google Borg traces~\cite{Tirmazi20,Liu12} are the formal dataset; this artefact uses a 2011 time-sliced subset, not the full 2019 eight-cell dump.
